Geophysical inversion is underdetermined: observations constrain a
subsurface model only through acquisition geometry, a forward operator,
error assumptions, parameterization, prior information, and numerical
optimization. Inverse Earth Studio is an open-source research workbench
connecting data ingestion and processing, forward modelling, inversion,
diagnostics, and interpretation across potential fields,
electromagnetics, electrical methods, and seismology. Its implementation
includes Python pipelines and solvers, bilingual computational
notebooks, provenance-bearing experiments, browser instruments, and a
bounded authenticated-service design.

The method suite spans gravity and magnetic processing and inversion,
magnetotelluric screening and layered-earth inversion,
electrical-resistivity quality control and inversion, waveform
processing and phase picking, travel-time inversion, acoustic
full-waveform inversion, joint inversion, and learned estimators.
Results include both recoveries and failures. For example, a synthetic
gravity case has weighted WRMS 1 while recovered model correlation is
0.283 and centroid depth is about 570 m for a 282 m truth; a
low-velocity FWI case remains unresolved. Clear Lake station cl061 is
retained for tensor quality control and excluded from 1D inversion after
failing the consistency screen. A proposed travel-time residual
refinement shows a small positive result on one fixed synthetic cohort,
which does not establish field transfer or general superiority.

The workbench makes data fit, model recovery, and geological
interpretation separate questions. Synthetic cases carry known truth;
field observations do not. Replayed browser artifacts are distinguished
from online calculations, and notebook execution is distinguished from
scientific validation. The report documents the methods, numerical
assumptions, data provenance, executable results, and limitations; it
does not claim that unresolved scientific or service-acceptance gates
have passed.
